\documentclass[10pt,a4paper]{amsart}
\usepackage[margin=19mm]{geometry}
\usepackage{amsmath,amssymb,mathtools}
\usepackage[scr=rsfs,cal=euler]{mathalfa}
\usepackage{xfrac}
\usepackage[unicode,hidelinks,bookmarks=false]{hyperref}
\newcommand{\ie}{i.e.\ }
\newcommand{\cf}{cf.\ }
\newcommand{\resp}{respectively\ }
\newcommand{\Subsection}{\subsection}
\providecommand{\nobreakdash}{\nobreak\hbox{-}}
\setlength{\emergencystretch}{2em}
\multlinegap0pt
\usepackage{amssymb}
\newtheorem{proposition}{Proposition}
\newtheorem{lemma}{Lemma}
\newcommand{\C}{\mathbb C}
\title{Strongly linearly convex exhaustion of $\mathbb C$-convex domains with $C^{1,\alpha}$ boundary}
\author{Armen Edigarian}
\date{Source: arXiv:2607.26696v1 (CC BY 4.0)}
\begin{document}
\maketitle

\noindent\emph{Source and licence.} Strongly linearly convex exhaustion of $\mathbb C$-convex domains with $C^{1,\alpha}$ boundary. Version arXiv:2607.26696v1 (CC BY 4.0), distributed by arXiv under the Creative Commons Attribution 4.0 International licence, http://creativecommons.org/licenses/by/4.0/, as recorded in the arXiv licence metadata for this exact version and shown on the abstract page https://arxiv.org/abs/2607.26696v1. Full licence notice: \texttt{licenses/arxiv-2607.26696-CC-BY-4.0.txt}.\par\smallskip

\noindent\textit{Editorial scope.} Counted unit: the article's proposition prop:mollification (source0.tex lines 292--305). The article's complete original proof of it is delivered with it (source0.tex lines 307--425, one \texttt{proof} environment of 119 lines). No mathematics is written, repaired or re-derived: the statement and the proof are the article's own lines, in the article's own notation, and no retained line is substituted or re-flowed.  Retained uncounted as the source-local context the proof needs: lines 70--73; lines 133--141; lines 193--205; lines 258--260.  Nothing was omitted from any retained span, so the package carries no omission record.  No retained line contains a citation, so the wrapper carries no bibliography and no bibliography entry.  No retained line contains a figure environment, an image-inclusion command or drawing code, so no image is delivered and the package declares no asset dependency.  Editorial formatting only, in addition: an AMS \texttt{amsart} preamble that restates the article's own declarations and macros used by the retained lines, namely the environments \texttt{proposition}, \texttt{lemma} on separate counters, together with the standard packages those lines need.  The retained environments print in this document as Proposition 1, Lemma 1 in order of appearance, whereas the article prints the same items with the numbering of its own full document.  The complete native source of the article as distributed in the arXiv e-print for this exact version is bundled unchanged as \texttt{sources/206290/source0.tex}.
\begin{equation}\label{eq:complex-line-second}
 \frac{d^2}{dt^2}u(z+t e^{i\theta}w)\bigg|_{t=0}
 =2L_u(z;w)+2\operatorname{Re}\bigl(e^{2i\theta}Q_u(z;w)\bigr).
\end{equation}

\begin{proposition}\label{prop:second-difference}
Let $z\in D$ be a differentiability point of $h$. Then, for any
$\theta\in\mathbb R$, any $w\in\C^n$, and any real $t$ sufficiently small
that $z\pm te^{i\theta}w\in D$,
\begin{equation}\label{eq:second-difference}
 h(z+t e^{i\theta}w)+h(z-t e^{i\theta}w)-2h(z)
 \le 2t^2\frac{|\partial h(z)(w)|^2}{h(z)}.
\end{equation}
\end{proposition}

\begin{lemma}\label{lem:gradosc}
There is a constant $C>0$ such that, for all sufficiently small $\delta>0$,
whenever $z,\zeta$ are differentiability points of $h=\delta_D^2$ satisfying
\[
 \frac\delta2\le\delta_D(z),\delta_D(\zeta)\le5\delta,
 \qquad |z-\zeta|\le\frac\delta4,
\]
then, for every $w\in\C^n$,
\begin{equation}\label{eq:gradestimate}
 |\partial h(z)(w)-\partial h(\zeta)(w)|
 \le C\bigl(|z-\zeta|+\delta^\gamma\bigr)|w|.
\end{equation}
\end{lemma}

\begin{equation}\label{eq:hosc}
 |h(z)-h(\zeta)|\le10\delta|z-\zeta|
\end{equation}

\begin{proposition}\label{prop:mollification}
There is a constant $C>0$ such that, for all sufficiently small $\delta$,
every $z\in U_\delta$ and every $w\in\C^n$ satisfy
\begin{equation}\label{eq:mollifiedineq}
 \frac{|\partial h_\delta(z)(w)|^2}{h_\delta(z)}
 -L_{h_\delta}(z;w)-|Q_{h_\delta}(z;w)|
 \ge -C\delta^{\frac{2\alpha}{1-\alpha}}|w|^2.
\end{equation}
Moreover
\begin{equation}\label{eq:C0approx}
 |h_\delta(z)-h(z)|\le10\delta\tau_\delta,
 \qquad z\in U_\delta.
\end{equation}
\end{proposition}

\begin{proof}
Fix $z\in U_\delta$, $w\in\C^n$, and $\theta\in\mathbb R$. Since
$\delta_D(\zeta)\ge\delta-\tau_\delta$ for every
$\zeta\in B(z,\tau_\delta)$, there is $t_0>0$, depending on
$z,\delta,w$ but not on $\zeta$, such that
$\zeta\pm te^{i\theta}w\in D$ for every
$\zeta\in B(z,\tau_\delta)$ whenever $|t|<t_0$. For such $t$,
Proposition~\ref{prop:second-difference} applies at almost every
$\zeta\in B(z,\tau_\delta)$. Averaging with
$\chi_{\tau_\delta}(z-\zeta)$ gives
\begin{align*}
 &h_\delta(z+t e^{i\theta}w)+h_\delta(z-t e^{i\theta}w)-2h_\delta(z)\\
 &\qquad\le
 2t^2\left(\frac{|\partial h(w)|^2}{h}*\chi_{\tau_\delta}\right)(z).
\end{align*}
Divide by $t^2$, let $t\to0$, and use
\eqref{eq:complex-line-second}. We obtain
\[
 L_{h_\delta}(z;w)+\operatorname{Re}\bigl(e^{2i\theta}Q_{h_\delta}(z;w)\bigr)
 \le\left(\frac{|\partial h(w)|^2}{h}*\chi_{\tau_\delta}\right)(z).
\]
Taking the maximum over $\theta$ yields
\begin{equation}\label{eq:convineq}
 L_{h_\delta}(z;w)+|Q_{h_\delta}(z;w)|
 \le\left(\frac{|\partial h(w)|^2}{h}*\chi_{\tau_\delta}\right)(z).
\end{equation}

It remains to compare the right-hand side with
$|\partial h_\delta(z)(w)|^2/h_\delta(z)$. Expanding the square gives the
exact identity
\begin{align}\label{eq:variance-identity}
 &\int
 \frac{|\partial h(\zeta)(w)|^2}{h(\zeta)}
 \chi_{\tau_\delta}(z-\zeta)\,d\zeta
 -\frac{|\partial h_\delta(z)(w)|^2}{h_\delta(z)}\notag\\
 &\qquad=
 \int h(\zeta)
 \left|
 \frac{\partial h(\zeta)(w)}{h(\zeta)}
 -\frac{\partial h_\delta(z)(w)}{h_\delta(z)}
 \right|^2
 \chi_{\tau_\delta}(z-\zeta)\,d\zeta.
\end{align}

If $\zeta\in B(z,\tau_\delta)$, then for small $\delta$
\[
 \frac{3\delta}{4}<\delta_D(\zeta)<\frac{17\delta}{4},
\]
so
\begin{equation}\label{eq:hexplicitbounds}
 \frac{9}{16}\delta^2<h(\zeta)<\frac{289}{16}\delta^2
\end{equation}
and, at differentiability points,
\[
 |\partial h(\zeta)(w)|\le\frac{17}{4}\delta|w|.
\]
For $\zeta,\eta\in B(z,\tau_\delta)$ we have
$|\zeta-\eta|\le2\tau_\delta$. For sufficiently small $\delta$,
$2\tau_\delta\le\delta/4$, so Lemma~\ref{lem:gradosc} gives
\[
 |\partial h(\zeta)(w)-\partial h(\eta)(w)|
 \le C\delta^\gamma|w|
\]
for almost every pair $\zeta,\eta$. Also
\eqref{eq:hosc} gives
\[
 |h(\zeta)-h(\eta)|\le C\delta\tau_\delta.
\]
Consequently
\begin{align*}
 \left|
 \frac{\partial h(\zeta)(w)}{h(\zeta)}
 -\frac{\partial h(\eta)(w)}{h(\eta)}
 \right|
 &\le
 \frac{|\partial h(\zeta)(w)-\partial h(\eta)(w)|}{h(\zeta)}\\
 &\quad+
 \frac{|\partial h(\eta)(w)|\,|h(\zeta)-h(\eta)|}
 {h(\zeta)h(\eta)}\\
 &\le C\bigl(\delta^{\gamma-2}+\tau_\delta\delta^{-2}\bigr)|w|\\
 &\le C\delta^{\gamma-2}|w|.
\end{align*}
Now
\[
 \frac{\partial h_\delta(z)(w)}{h_\delta(z)}
 =\frac{1}{h_\delta(z)}
 \int h(\eta)\frac{\partial h(\eta)(w)}{h(\eta)}
 \chi_{\tau_\delta}(z-\eta)\,d\eta,
\]
so it is the average of $\partial h(\eta)(w)/h(\eta)$ with respect to the
probability measure
\[
 \frac{h(\eta)\chi_{\tau_\delta}(z-\eta)}{h_\delta(z)}\,d\eta.
\]
Therefore the same oscillation bound gives
\[
 \left|
 \frac{\partial h(\zeta)(w)}{h(\zeta)}
 -\frac{\partial h_\delta(z)(w)}{h_\delta(z)}
 \right|
 \le C\delta^{\gamma-2}|w|
\]
for almost every $\zeta\in B(z,\tau_\delta)$. Using
\eqref{eq:hexplicitbounds} in \eqref{eq:variance-identity}, we obtain
\[
 \left(\frac{|\partial h(w)|^2}{h}*\chi_{\tau_\delta}\right)(z)
 -\frac{|\partial h_\delta(z)(w)|^2}{h_\delta(z)}
 \le C\delta^{2\gamma-2}|w|^2.
\]
Combining this with \eqref{eq:convineq}, and using
$2\gamma-2=\frac{2\alpha}{1-\alpha}$, proves
\eqref{eq:mollifiedineq}.

Finally, if $\zeta\in B(z,\tau_\delta)$, then \eqref{eq:hosc} gives
\[
 |h(\zeta)-h(z)|\le10\delta|\zeta-z|\le10\delta\tau_\delta.
\]
Averaging proves \eqref{eq:C0approx}.
\end{proof}

\end{document}
